\documentclass[11pt]{article}
\usepackage[margin=25mm]{geometry}
\usepackage{fontspec}
\setmainfont{Latin Modern Roman}
\usepackage{amsmath,amssymb,mathtools}
\usepackage{unicode-math}
\setmathfont{Latin Modern Math}
\usepackage{xurl,hyperref}
\hypersetup{colorlinks=true,urlcolor=blue}
\setcounter{secnumdepth}{0}
\setlength{\parindent}{0pt}
\setlength{\parskip}{5pt}
\setlength{\emergencystretch}{3em}
\newcommand{\sref}[2]{\hyperlink{source:#1:#2}{[S#2]}}
\newcommand{\cataloguescope}[1]{\paragraph{Recorded target.}#1}
\newcommand{\R}{\mathbb R}
\newcommand{\C}{\mathbb C}
\newcommand{\N}{\mathbb N}
\newcommand{\Z}{\mathbb Z}
\newcommand{\Q}{\mathbb Q}
\newcommand{\F}{\mathbb F}
\newcommand{\eps}{\varepsilon}
\begin{document}
\section{Matrix virial decomposition and a collision-free tensor stress test}\label{statement}

\subsection{Definitions and mathematical statement}

\paragraph{Shared notation.}$\mathbb N=\{1,2,\ldots\}$; the standard number fields have their usual meanings. All additional conventions are stated locally.


\paragraph{Definitions and notation.}In the Newtonian $n$-body problem, positions $q_i\in\mathbb R^d$ and masses $m_i>0$ obey $m_i\ddot q_i=\nabla_{q_i}U$, where $U(q)=\sum_{i<j}m_im_j/|q_i-q_j|$ and the gravitational constant is normalized to one. Here $T=\tfrac12\sum_i m_i|\dot q_i|^2$, $E=T-U$, and a collision has $q_i=q_j$. The center of mass can be fixed at zero; moment of inertia is $I=\sum_i m_i|q_i|^2$. A central configuration has $\nabla_{q_i}U=-\lambda m_iq_i$ in center-of-mass coordinates for a common scalar $\lambda$. After fixing $I=1$, it is a critical point of $U$ on the collision-free inertia ellipsoid. Congruent configurations and scales are factored out for a finiteness count. A relative equilibrium is a rigidly rotating solution with time-independent shape. Saari’s conjecture in its moment-of-inertia formulation asks whether constant scalar $I$ along a Newtonian solution forces a relative equilibrium. The inertia tensor in center-of-mass coordinates is the matrix $Q=\sum_i m_iq_iq_i^{\mathsf T}$ (or its usual trace-complement physical convention); fixing it is a distinct, stronger geometric condition. A tensor convention matters in higher dimension.

\paragraph{Catalogue wording retained for traceability.}
Give a conceptual proof of Saari’s conjecture. Also fix the moment-of-inertia tensor and investigate the corresponding question, potentially in higher dimensions.
\paragraph{Scope and source qualification.}The source has both a scalar conjecture and a proposed tensor variant. A tensor result is not a proof of the scalar conjecture, and dimensional/mass hypotheses must be stated for any claimed theorem. The catalogue title/status was paired with a machine-generated research attempt; that attempt is not peer-reviewed evidence of settlement. The statement covers the original source questions rather than only the title’s subcase.
\subsection{Short English statement}
Find a conceptual proof of Saari’s conjecture and investigate the fixed-inertia-tensor analogue.
\subsection{Research attempt}
\paragraph{Model, date and investigated scope.}
OpenAI GPT-6 Astra (configured model \texttt{gpt-6-astra}), ultra reasoning; generated 11 October 2026. Derived the exact matrix virial identity and constant-I/constant-Q necessary consequences. No scalar Saari rigidity or tensor analogue is proved.

\paragraph{Sources and literature check.}
Read the supplied complete definition, catalogue wording, scope qualifications and source list for this canonical ID. The original-source citations are inherited as supplied references; no fresh primary-source or novelty search was performed for this attempt. No current open-status claim is inferred from this calculation.

\paragraph{The exact tensor virial identity.}For a collision-free Newtonian solution define \(Q=\sum_im_iq_iq_i^T\), \(K=\sum_im_i\dot q_i\dot q_i^T\), and
\[S=\sum_{i<j}\frac{m_im_j}{|q_i-q_j|^3}(q_i-q_j)(q_i-q_j)^T.
\]
Differentiating \(Q\) twice and grouping the force terms by pairs gives \(Q''=2K-2S\): the pair forces are opposite and their outer products combine to minus twice the displayed pair matrix. Taking traces gives \(I''=4T-2U\), since \(\operatorname{tr}K=2T\) and \(\operatorname{tr}S=U\).
\paragraph{Consequences and their exact limit.}If scalar \(I\) is constant, then \(2T=U\). Conservation of \(E=T-U\) implies \(U=-2E\) and \(T=-E\) are individually constant. If the entire tensor \(Q\) is constant, the stronger matrix equation \(K=S\) holds at every time. In particular \(S\) is positive definite whenever the mutual-difference vectors span the ambient space, because its quadratic form is a positive weighted sum of squares.
None of these instantaneous equalities forces all mutual distances to be constant. The trace discards directional information, and even \(K=S\) compares velocity covariance with force stress without identifying a common angular velocity. Treating this necessary balance as rigidity would be circular.
\paragraph{Gap.}The first missing step is a dynamical rigidity argument converting the constant-potential/kinetic constraints into fixed shape and uniform rotation, or determining precisely why the tensor variant is different in higher dimension. No general scalar Saari proof, tensor classification or mass/dimension extension is claimed. The identity is a standard starting point reconstructed with the source's tensor convention.

\paragraph{Verification and adversarial audit.}
Checked pairwise force grouping, factors of two, trace, positivity criterion and scalar versus tensor scope. This was an LLM adversarial audit, not independent human verification; no theorem prover was used. Prompts A then B applied to this ID: Checked pairwise force grouping, factors of two, trace, positivity criterion and scalar versus tensor scope. Scope remains limited: Convert constant energy/potential and tensor stress balance into shape rigidity using additional dynamical information.

Coordinator checked the tensor and scalar virial identities, including the positive-potential sign and fixed-inertia consequences; no rigidity implication follows. This additional check was model-assisted, not independent human verification.

\paragraph{Tools, process and contributions.}
Python 3 and local shell used to read supplied files and serialize this manuscript; mathematics derived symbolically in the text, with no numerical experiment or formal proof assistant. Contribution prompt A was followed by prompt B. GPT-6 Astra with ultra reasoning drafted the research and performed its stated self-audit; any additional coordinator contribution is identified in the verification and audit. A GPT-6-based Codex coordinator assembled the manuscript and checked repository integration; its exact variant and reasoning configuration were not exposed. No human mathematical contribution, independent human proof check or formal proof certification is supplied. Complete derivations and the consulted source references are included here; private reasoning transcripts are not published.

\paragraph{Final status and first remaining gap.}
\textbf{UNRESOLVED}. Convert constant energy/potential and tensor stress balance into shape rigidity using additional dynamical information.

\section{Completion Estimate}
\noindent\textbf{COMPLETION: 0\%}
The calculation is an elementary benchmark or conditional reduction; it establishes no claimed new portion of the complete source target.

\subsection{Sources}
\begin{itemize}
\item[\hypertarget{source:statement:1}{S1}] ULAM.AI and the UnsolvedMath Contributors, UnsolvedMath, catalogue entry. Catalogue curation under CC BY 4.0; definitions and English exposition added locally.
\url{https://huggingface.co/datasets/ulamai/UnsolvedMath}.
\item[\hypertarget{source:statement:2}{S2}] American Institute of Mathematics, Variational Methods in Celestial Mechanics; notes by Kuo-Chang Chen, version 22 June 2003, source item 14.
\url{https://aimath.org/WWN/varcelest/varcelest.pdf}.
\end{itemize}
\label{end:statement}
\end{document}
